% !TEX program = lualatex
% ---------------------------------------------------------------------------------------------
%  Readback note 2: The four statements: the Riemann hypothesis and its generalisation to Dirichlet L-functions
%
%  Build:  python3 tools/rb.py note <slug>
%
%  One unit of a Lean artifact, rewritten as ordinary mathematics. The prose is hand-written. Lean names,
%  line ranges and counts come from data.tex; quoted Lean statements are copied from the artifact by the build.
%
%  Rules of the text: standard terminology; every symbol is explained where it first appears and listed again
%  in the notation at the end; no letter has two meanings; a proof is complete or says what it takes for granted;
%  Lean names appear only in the LEAN lines and in the concordance.
%
%  The signature rule (the build enforces it): a function, relation or predicate that the artifact defines is
%  written only through its notation macro, which the build makes from the spec (`notation`) with one slot for
%  each explicit argument of the Lean definition. No argument is fixed and dropped, also not "for brevity" and
%  not for one section. Every argument and every hypothesis of a stated declaration stands inside its
%  statement, or in a Setting that the statement cites by number; only implicit and instance arguments (the
%  ambient objects) may stand in the conventions.
% ---------------------------------------------------------------------------------------------
\documentclass[11pt,a4paper]{article}
\ifdefined\pdfvariable\pdfvariable suppressoptionalinfo \numexpr 1+2+4+8\relax\fi
\usepackage[a4paper,hmargin=35mm,top=22mm,bottom=27mm,footskip=12mm,nohead]{geometry}

\usepackage{fontspec}
\directlua{
  luaotfload.add_fallback("gfallback", {"STIXTwoMath-Regular.otf:mode=node;", "DejaVuSans.ttf:mode=node;"})
  luaotfload.add_fallback("gmonofallback", {"DejaVuSansMono.ttf:mode=node;", "STIXTwoMath-Regular.otf:mode=node;"})
}
\setmainfont{NewCM10}[Extension=.otf,UprightFont=*-Book,ItalicFont=*-BookItalic,
  BoldFont=*-Bold,BoldItalicFont=*-BoldItalic,SmallCapsFont=*-Book,SmallCapsFeatures={Letters=SmallCaps},
  RawFeature={fallback=gfallback}]
\setsansfont{NewCMSans10}[Extension=.otf,UprightFont=*-Book,ItalicFont=*-BookOblique,
  BoldFont=*-Bold,BoldItalicFont=*-BoldOblique,RawFeature={fallback=gfallback}]
\setmonofont{JetBrainsMono}[Extension=.otf,UprightFont=*-Regular,ItalicFont=*-Italic,
  BoldFont=*-SemiBold,BoldItalicFont=*-SemiBoldItalic,Scale=MatchLowercase,
  RawFeature={-calt;-liga;fallback=gmonofallback}]
\usepackage{amsmath}
\usepackage{amsthm}
\usepackage[math-style=TeX]{unicode-math}
\setmathfont{NewCMMath-Book.otf}
\emergencystretch=2.5em\relax

\usepackage{xcolor}
\usepackage{tikz}
\usetikzlibrary{calc,arrows.meta,positioning,backgrounds}
\usepackage[protrusion=true,expansion=true]{microtype}
\usepackage{enumitem}
\usepackage{needspace}
\usepackage{booktabs}
\usepackage{array}
\usepackage{longtable}
\usepackage{fancyhdr}
\usepackage{lastpage}
\usepackage{multicol}
\usepackage[hidelinks,pdfusetitle]{hyperref}

\definecolor{ink}{HTML}{0F172A}
\definecolor{inksoft}{HTML}{475569}
\definecolor{inkfaint}{HTML}{94A3B8}
\definecolor{rulegrey}{HTML}{CBD5E1}
\definecolor{arith}{HTML}{C2410C}
\definecolor{est}{HTML}{7C3AED}
\definecolor{ana}{HTML}{0891B2}
\definecolor{ok}{HTML}{047857}
\definecolor{leancol}{HTML}{3730A3}

\title{The four statements: the Riemann hypothesis and its generalisation to Dirichlet L-functions}
\author{grh-compact}
\date{}

% ------------------------------------------------------------------------------------- macros
\newcommand{\leanfont}{\ttfamily\fontsize{8.6}{10.5}\selectfont\color{leancol}}
\newcommand{\lean}[1]{\csname LN@#1\endcsname}
\newcommand{\fillin}[1]{{\color{arith}\itshape[#1]}}  % a place of the template that is not written yet; the build reports every one that is left
\newcommand{\leanquote}[1]{\csname LQ@#1\endcsname}
% generated by rb.py note; do not edit
\newcommand{\ConcDecls}{4}
\newcommand{\Pin}{72fead10}
\newcommand{\PinText}{that of commit {\leanfont 72fead10}}
\newcommand{\ArtifactName}{grh-compact}
\expandafter\newcommand\csname LN@rh\endcsname{\href{https://github.com/lsernau/grh-lean/blob/72fead102f663d5c7a407964a13dc0cf9d977207/grh-compact/src/Statements.lean\#L3-L4}{{\leanfont RiemannZetaZerosOnLine}}}
\expandafter\newcommand\csname LQ@rh\endcsname{\leanquoteline{0}{def RiemannZetaZerosOnLine : Prop :=}\leanquoteline{2}{∀ s : ℂ, riemannZeta s = 0 → (¬∃ n : ℕ, s = -2 * ((n : ℂ) + 1)) → s ≠ 1 → s.re = 1 / 2}}
\expandafter\newcommand\csname LN@halfPlane\endcsname{\href{https://github.com/lsernau/grh-lean/blob/72fead102f663d5c7a407964a13dc0cf9d977207/grh-compact/src/Statements.lean\#L6-L8}{{\leanfont DirichletLNonvanishingRightOfHalf}}}
\expandafter\newcommand\csname LQ@halfPlane\endcsname{\leanquoteline{0}{def DirichletLNonvanishingRightOfHalf : Prop :=}\leanquoteline{2}{∀ (N : ℕ) [NeZero N] (χ : DirichletCharacter ℂ N) (s : ℂ),}\leanquoteline{4}{1 / 2 < s.re → (χ ≠ 1 ∨ s ≠ 1) → DirichletCharacter.LFunction χ s ≠ 0}}
\expandafter\newcommand\csname LN@strip\endcsname{\href{https://github.com/lsernau/grh-lean/blob/72fead102f663d5c7a407964a13dc0cf9d977207/grh-compact/src/Statements.lean\#L10-L12}{{\leanfont DirichletLZerosOnLine}}}
\expandafter\newcommand\csname LQ@strip\endcsname{\leanquoteline{0}{def DirichletLZerosOnLine : Prop :=}\leanquoteline{2}{∀ (N : ℕ) [NeZero N] (χ : DirichletCharacter ℂ N) (s : ℂ),}\leanquoteline{4}{0 < s.re → s.re < 1 → DirichletCharacter.LFunction χ s = 0 → s.re = 1 / 2}}
\expandafter\newcommand\csname LN@stripPrim\endcsname{\href{https://github.com/lsernau/grh-lean/blob/72fead102f663d5c7a407964a13dc0cf9d977207/grh-compact/src/Statements.lean\#L14-L16}{{\leanfont PrimitiveDirichletLZerosOnLine}}}
\expandafter\newcommand\csname LQ@stripPrim\endcsname{\leanquoteline{0}{def PrimitiveDirichletLZerosOnLine : Prop :=}\leanquoteline{2}{∀ (N : ℕ) [NeZero N] (χ : DirichletCharacter ℂ N), χ.IsPrimitive →}\leanquoteline{4}{∀ s : ℂ, 0 < s.re → s.re < 1 → DirichletCharacter.LFunction χ s = 0 → s.re = 1 / 2}}
\newcommand{\ConcordanceRows}{\crow{rh}{\href{https://github.com/lsernau/grh-lean/blob/72fead102f663d5c7a407964a13dc0cf9d977207/grh-compact/src/Statements.lean\#L3-L4}{{\leanfont RiemannZetaZerosOnLine}}}{Statements.lean}{3--4}
\crow{halfPlane}{\href{https://github.com/lsernau/grh-lean/blob/72fead102f663d5c7a407964a13dc0cf9d977207/grh-compact/src/Statements.lean\#L6-L8}{{\leanfont DirichletLNonvanishingRightOfHalf}}}{Statements.lean}{6--8}
\crow{strip}{\href{https://github.com/lsernau/grh-lean/blob/72fead102f663d5c7a407964a13dc0cf9d977207/grh-compact/src/Statements.lean\#L10-L12}{{\leanfont DirichletLZerosOnLine}}}{Statements.lean}{10--12}
\crow{stripPrim}{\href{https://github.com/lsernau/grh-lean/blob/72fead102f663d5c7a407964a13dc0cf9d977207/grh-compact/src/Statements.lean\#L14-L16}{{\leanfont PrimitiveDirichletLZerosOnLine}}}{Statements.lean}{14--16}}
% the notations of the defined functions: one slot for each explicit argument of the Lean declaration
\newcommand{\UnitMembers}{4}
\newcommand{\UnitInNote}{4}
\newcommand{\UnitProofs}{0}
\newcommand{\UnitSteps}{0}
\newcommand{\UnitAllSteps}{0}


\linespread{1.08}                       % 11 pt on 14.7 pt
\setlength{\parindent}{0pt}
\setlength{\parskip}{5.5pt plus 1.5pt minus 0.5pt}
\setlength{\abovedisplayskip}{7.5pt plus 2pt minus 3pt}
\setlength{\belowdisplayskip}{7.5pt plus 2pt minus 3pt}
\setlength{\abovedisplayshortskip}{3pt plus 2pt}
\setlength{\belowdisplayshortskip}{6pt plus 2pt minus 2pt}
\emergencystretch=1.5em
\clubpenalty=10000 \widowpenalty=10000 \displaywidowpenalty=10000
\relpenalty=7000 \binoppenalty=9000
\raggedbottom
\pagestyle{fancy}\fancyhf{}\renewcommand{\headrulewidth}{0pt}
\fancyfoot[C]{\sffamily\fontsize{7.6}{9}\selectfont\color{inkfaint}\thepage\,/\,\pageref*{LastPage}}
\setlist{topsep=2pt,partopsep=0pt,parsep=0pt,itemsep=4pt,leftmargin=2.1em,labelsep=0.55em}

\newcommand{\Q}{\mathbb{Q}}
\newcommand{\Z}{\mathbb{Z}}
\newcommand{\F}{\mathbb{F}}
\newcommand{\len}{\operatorname{len}}
\newcommand{\cok}{\operatorname{cok}}
\newcommand{\End}{\operatorname{End}}
\newcommand{\Hom}{\operatorname{Hom}}
\newcommand{\Gal}{\operatorname{Gal}}
\newcommand{\diag}{\operatorname{diag}}
\newcommand{\kw}[1]{\emph{#1}}
\newcommand{\m}[1]{\mbox{$#1$}}        % a short formula that must not be broken across lines
\newcommand{\head}[1]{\par\addvspace{17pt plus 4pt minus 2pt}\Needspace{6\baselineskip}%
  \noindent{\fontsize{13.2}{16}\selectfont\bfseries\color{ink}#1}\par\nopagebreak\vspace{-1pt}}
\newcommand{\sect}[2]{\refstepcounter{section}\label{#2}\par\addvspace{19pt plus 4pt minus 2pt}\Needspace{7\baselineskip}%
  \noindent\llap{{\fontsize{13.2}{16}\selectfont\bfseries\color{inkfaint}\thesection}\hspace{9pt}}%
  {\fontsize{13.2}{16}\selectfont\bfseries\color{ink}#1}\par\nopagebreak\vspace{-1pt}}
\newcommand{\subhead}[1]{\par\addvspace{9pt plus 2pt}\Needspace{4\baselineskip}\noindent\textbf{#1}\enspace\ignorespaces}
% where a statement lives in the Lean archive
\newcommand{\leantag}{{\sffamily\bfseries\fontsize{6.3}{8}\selectfont\color{leancol}\addfontfeatures{LetterSpace=7}LEAN}}
\newcommand{\sep}{\unskip\nobreak\hspace{5pt}{\color{inkfaint}\textperiodcentered}\hspace{5pt plus 2pt}\ignorespaces}
\newcommand{\leanline}[1]{\par\nopagebreak\vspace{-2.5pt}{\fontsize{8.9}{12.4}\selectfont\color{inksoft}\raggedright
  \noindent\llap{\leantag\hspace{9pt}}#1\par}\vspace{1pt}}
% a displayed statement in a tinted box: #1 colour, #2 text
\newcommand{\statement}[2]{\par\addvspace{7pt}\begingroup\setlength{\fboxsep}{8.5pt}\setlength{\fboxrule}{0.6pt}%
  \noindent\fcolorbox{#1!30!white}{#1!3!white}{\parbox{\dimexpr\linewidth-2\fboxsep-2\fboxrule\relax}{#2}}%
  \endgroup\par\addvspace{7pt}}
\newcommand{\nt}[2]{\par\noindent\hangindent=1.1em\hangafter=1 #1\enspace{\color{inksoft}#2}\par\vspace{1pt}}
\newcommand{\legdot}[1]{\tikz[baseline=-0.62ex]\fill[#1] (0,0) circle (2.5pt);\hspace{3pt}}
\newcommand{\crow}[4]{\csname CS@#1\endcsname & #2 & \textcolor{inksoft}{#3} & \textcolor{inksoft}{#4}\\}
\newcommand{\leanquoteline}[2]{\par\noindent\hangindent=\dimexpr#1\fontdimen2\font+2.2em\relax\hangafter=1
  \hspace*{#1\fontdimen2\font}#2}
\newcommand{\quoted}[1]{\par\vspace{3pt}{\ttfamily\fontsize{7.5}{10.4}\selectfont\color{leancol}\raggedright
  \setlength{\parskip}{0pt}\leanquote{#1}\par}\vspace{4pt}}
\newcommand{\bx}[3]{{\sffamily\bfseries\fontsize{7.3}{9.2}\selectfont\color{#1!82!black}#2}\\[1.5pt]#3}
\newcommand{\cs}[2]{\expandafter\newcommand\csname CS@#1\endcsname{#2}}

% theorems
\newtheoremstyle{plainx}{7pt plus 2pt}{7pt plus 2pt}{\itshape}{}{\bfseries}{.}{0.5em}{}
\newtheoremstyle{uprightx}{7pt plus 2pt}{7pt plus 2pt}{}{}{\bfseries}{.}{0.5em}{}
\theoremstyle{plainx}
\newtheorem{theorem}{Theorem}[section]
\newtheorem{lemma}[theorem]{Lemma}
\newtheorem{proposition}[theorem]{Proposition}
\newtheorem{corollary}[theorem]{Corollary}
\theoremstyle{uprightx}
\newtheorem{definition}[theorem]{Definition}
\newtheorem{setting}[theorem]{Setting}          % hypotheses shared by several statements; each of them cites it by number
\newtheorem{remark}[theorem]{Remark}
\newtheorem{example}[theorem]{Example}
\renewcommand{\qedsymbol}{{\setlength{\fboxsep}{0pt}\setlength{\fboxrule}{0.5pt}\fbox{\rule{0pt}{0.46em}\hspace{0.46em}}}}
% what holds throughout the note: the ambient objects only (implicit and instance arguments), never a hypothesis
\newenvironment{conventions}{\par\addvspace{5pt}\noindent\textbf{Conventions.}\enspace\ignorespaces}{\par\addvspace{3pt}}
\makeatletter
\renewenvironment{proof}[1][Proof]{\par\pushQED{\qed}\normalfont\topsep0pt\partopsep0pt
  \trivlist\item[\hskip\labelsep\itshape #1.]\ignorespaces}{\popQED\endtrivlist\@endpefalse}
\makeatother

% where each declaration of the concordance is treated: \cs{key}{place}, the place naming a \label or a lettered item
\cs{rh}{Statement 1}
\cs{halfPlane}{Statement 2}
\cs{strip}{Statement 3}
\cs{stripPrim}{Statement 4}
\newcommand{\CC}{\mathbb{C}}
\newcommand{\ZZ}{\mathbb{Z}}
\newcommand{\re}{\operatorname{Re}}
\newcommand{\mlib}[1]{\leavevmode{\leanfont #1}}
\newcommand{\mrow}[2]{#1 & \mlib{#2}\\[1.5pt]}

\begin{document}
\color{ink}

{\raggedright
{\sffamily\bfseries\fontsize{7.2}{9}\selectfont\color{leancol}\addfontfeatures{LetterSpace=8}READBACK NOTE 2\par}
\vspace{5pt}
{\fontsize{21.5}{25.5}\selectfont\bfseries The four statements: the Riemann hypothesis and its generalisation to
Dirichlet $L$-functions\par}
\vspace{5pt}
{\fontsize{12.4}{16}\selectfont\itshape\color{inksoft} One unit of the Lean artifact \ArtifactName, read back into
mathematics\par}}
\vspace{8pt}
{\color{rulegrey}\hrule height 0.6pt}
\vspace{7pt}
\thispagestyle{empty}

\textbf{Summary.} The headline of the artifact is five theorems. Their statements are four propositions about the
Riemann zeta function and the Dirichlet $L$-functions, written in one short file that uses nothing but the
mathematical library of Lean: the Riemann hypothesis for~$\zeta$; the non-vanishing of every Dirichlet $L$-function
on the half-plane \m{\re s>\tfrac12}; and the statement that the zeros of these functions in the critical strip lie
on the line \m{\re s=\tfrac12}, for all characters and for primitive ones. This note states the four propositions in
ordinary language, says which objects of the library they speak of and what determines those objects, and records
how the four are related: the second implies the other three.

\statement{leancol}{\fontsize{9.6}{13.2}\selectfont\textbf{What this note is.} The artifact \ArtifactName{} is a
body of Lean code: the compact edition of a formal proof, due to L.~Alpöge with the AI system Claude, whose final
theorems are the Riemann hypothesis and its generalisation to Dirichlet $L$-functions. That headline is a claim. The
authors report that every file compiles with Lean v4.35.0 against a pinned version of the library Mathlib and that
Lean's kernel accepted the proofs; an independent build is recorded in the handout's record of checks when it
reports. No mathematician has read the argument through, and nothing here says that the Riemann hypothesis is
proved. This note treats one unit of the artifact: all \UnitInNote{} of its \UnitMembers{} declarations, which are
definitions of propositions. The unit proves nothing, so there is no proof to read and nothing is taken for granted
from other units. Each of the four statements was compared with its Lean text, and the objects of the library that
they use were compared with the library's source at the pinned commit (Appendix~A says exactly what was read). The
Lean text read is \PinText; Appendix~A is the concordance of \ConcDecls{} declarations.}

\sect{The four statements}{s:intro}
Throughout, $\zeta$ is the Riemann zeta function and $L(s,\chi)$ is the $L$-function of a Dirichlet
character~$\chi$, as the library defines them; §\ref{s:setting} says what these functions are. The first statement
is about~$\zeta$, the other three are about the $L$-functions.

\statement{arith}{\phantomsection\label{st:1}\textbf{Statement 1} (the Riemann hypothesis for $\zeta$).\enspace
\emph{For every complex number~$s$: if \m{\zeta(s)=0}, if there is no natural number~$n$ with \m{s=-2(n+1)}, and if
\m{s\ne1}, then \m{\re s=\tfrac12}.}}
\leanline{\lean{rh}}

\statement{arith}{\phantomsection\label{st:2}\textbf{Statement 2} (the half-plane statement).\enspace
\emph{For every integer \m{M\ge1}, every Dirichlet character~$\chi$ modulo~$M$ and every complex number~$s$ with
\m{\re s>\tfrac12}: if $\chi$ is not the principal character modulo~$M$ or \m{s\ne1}, then \m{L(s,\chi)\ne0}.}

\vspace{4pt}\phantomsection\label{st:3}\textbf{Statement 3} (the strip statement).\enspace
\emph{For every integer \m{M\ge1}, every Dirichlet character~$\chi$ modulo~$M$ and every complex number~$s$: if
\m{0<\re s}, \m{\re s<1} and \m{L(s,\chi)=0}, then \m{\re s=\tfrac12}.}

\vspace{4pt}\phantomsection\label{st:4}\textbf{Statement 4} (the strip statement for primitive characters).\enspace
\emph{For every integer \m{M\ge1} and every Dirichlet character~$\chi$ modulo~$M$ that is primitive, and for every
complex number~$s$: if \m{0<\re s}, \m{\re s<1} and \m{L(s,\chi)=0}, then \m{\re s=\tfrac12}.}}
\leanline{\lean{halfPlane}\sep \lean{strip}\sep \lean{stripPrim}}

Each of the four is a proposition, not a theorem: the unit defines them and asserts nothing. They are asserted by
the five final theorems of the artifact, which are the subject of the note \emph{The five final theorems, and what
one must trust}. The argument is for Statement~2; its last steps, and the passage from Statement~2 to Statements~1
and~3, are the subject of the note \emph{The main inequality, the bound for the primes, and the zero-free
half-plane}.

In classical terms, Statement~1 is the Riemann hypothesis, and Statement~3 is the generalised Riemann hypothesis
for Dirichlet $L$-functions in its usual form, for the zeros in the open critical strip. Statement~2 is the form in
which the artifact proves it: a zero-free half-plane. It contains the classical non-vanishing on the closed
half-plane \m{\re s\ge1}, and beyond that it says that there is no zero with \m{\tfrac12<\re s<1}.

\sect{The objects of the statements}{s:setting}
\begin{conventions}\phantomsection\label{c:conventions}
The natural numbers are \m{0,1,2,\dots}, and $\re s$ is the real part of a complex number~$s$. The declarations of
this unit have no arguments, so there is no ambient object to fix. Everything that the four statements speak of is
defined in Mathlib, the mathematical library of Lean, at the commit that the artifact pins; nothing in them is
defined by the artifact. Appendix~A gives the library's names for the objects and facts of this section.
\end{conventions}

\subhead{The zeta function.} The library's zeta function is a function \m{\zeta\colon\CC\to\CC}, defined at
every complex number. It is constructed from Hurwitz zeta functions, and the library proves the following about it.
\begin{enumerate}[label=\textup{(\alph*)},leftmargin=2.3em,itemsep=2pt]
\item For \m{\re s>1}, \m{\zeta(s)=\sum_{n\ge1}n^{-s}}.
\item $\zeta$ is complex differentiable at every \m{s\ne1}, and \m{(s-1)\zeta(s)\to1} as \m{s\to1}.
\item \m{\zeta(-2(n+1))=0} for every natural number~$n$, and \m{\zeta(0)=-\tfrac12}.
\end{enumerate}
By (a), (b) and the identity theorem, the restriction of $\zeta$ to \m{\CC\smallsetminus\{1\}} is \emph{the}
analytic continuation of the series: the two properties determine it. A function of Lean must have a value
everywhere, so the library also assigns a value to~$\zeta$ at the pole \m{s=1}; that value is a convention, and no
statement of this note depends on it, because Statement~1 excludes \m{s=1}. The zeros in~(c) are the trivial zeros
\m{-2,-4,-6,\dots}; they are what the second hypothesis of Statement~1 excludes.

\subhead{Dirichlet characters.} Let \m{M\ge1} be an integer. A \kw{Dirichlet character modulo~$M$} (with complex
values) is a map \m{\chi\colon\ZZ/M\to\CC} with \m{\chi(1)=1} and \m{\chi(uv)=\chi(u)\chi(v)} for all classes $u$
and~$v$, and with \m{\chi(u)=0} for every class~$u$ that is not a unit. It is evaluated at an integer through the
class of the integer. The \kw{principal character} modulo~$M$ takes the value~$1$ at every unit; in the library it
is the unit element~$1$ of the group of characters. For \m{M=1} the ring \m{\ZZ/1} has one element, which is a
unit, and the principal character is the only character. If $M_0$ divides~$M$, a character modulo~$M_0$
\kw{induces} a character modulo~$M$, by composition with the reduction of units. The \kw{conductor} of a
character~$\chi$ modulo~$M$ is the least divisor~$M_0$ of~$M$ such that $\chi$ is induced by a character
modulo~$M_0$, and $\chi$ is \kw{primitive} when its conductor is~$M$. Every character modulo~$M$ is induced by
exactly one primitive character, whose modulus is the conductor.

\subhead{Dirichlet $L$-functions.} For an integer \m{M\ge1} and a Dirichlet character $\chi$ modulo~$M$, the
library's $L$-function is a function \m{s\mapsto L(s,\chi)} from $\CC$ to~$\CC$, defined at every complex number.
It is constructed as a linear combination of Hurwitz zeta functions, and the library proves the following.
\begin{enumerate}[label=\textup{(\alph*)},leftmargin=2.3em,itemsep=2pt]
\item For \m{\re s>1}, \m{L(s,\chi)=\sum_{n\ge1}\chi(n)\,n^{-s}}.
\item \m{s\mapsto L(s,\chi)} is complex differentiable at every~$s$, except possibly at \m{s=1} when $\chi$ is
principal.
\item For \m{M=1}, \m{L(s,\chi)=\zeta(s)} for every complex number~$s$.
\item If $\chi_0$ is a character modulo a divisor $M_0$ of~$M$ and $\chi$ is the character modulo~$M$ that it
induces, then for every complex number~$s$ such that $\chi_0$ is not principal or \m{s\ne1},
\[
L(s,\chi)=L(s,\chi_0)\prod_{p\mid M}\bigl(1-\chi_0(p)\,p^{-s}\bigr),
\]
the product over the primes~$p$ that divide~$M$.
\end{enumerate}
So for a character that is not principal, \m{L(\,\cdot\,,\chi)} is the entire function that continues the series.
For a principal character it is the continuation to \m{\CC\smallsetminus\{1\}}, with a simple pole at~$1$, and its
value at \m{s=1} is again a convention; this is why Statement~2 excepts the point \m{s=1} for a principal
character, and only that point.

\sect{How the four statements are related}{s:proofs}
The unit contains no proof. What follows is said for orientation; the proofs are in the notes named.

\subhead{Statement 2 implies the other three.} Statement~2 implies Statement~3: a zero~$s$ with
\m{0<\re s<\tfrac12} would be reflected by the functional equation to the zero \m{1-s} of the $L$-function of the
inverse character, which lies in the half-plane of Statement~2. Statement~2 for \m{M=1} implies Statement~1, by
property~(c) of the $L$-functions and the functional equation of~$\zeta$. These two implications are Theorems~4
and~5 of the note \emph{The main inequality, the bound for the primes, and the zero-free half-plane}. Statement~3
implies Statement~4, which is its special case for primitive characters. The artifact proves Statement~2 and
deduces the other three in this way.

\begin{remark}[what Statements 2 to 4 leave out]\label{r:left}
Statements 2, 3 and~4 say nothing about zeros with \m{\re s\le0}. There the $L$-function of a character has its
trivial zeros, and the $L$-function of a character that is not primitive can have further zeros on the line
\m{\re s=0}: those of the factors \m{1-\chi_0(p)p^{-s}} of property~(d), with $\chi_0$ the primitive character
that induces it, for the primes~$p$ that divide the modulus but not the conductor. For the other primes~$p$
dividing the modulus, \m{\chi_0(p)=0} and the factor is~$1$; so a character whose modulus and conductor have the
same prime divisors, for instance a character modulo~$9$ of conductor~$3$, has no further zero. Statement~1 does
speak of all zeros
of~$\zeta$: it says that a zero is a trivial zero or lies on the line \m{\re s=\tfrac12}.
\end{remark}

\begin{remark}[characters that are not primitive]\label{r:imprimitive}
The statements are made for all characters, primitive or not. This costs nothing and gains nothing. Let $\chi$ be a
character modulo~$M$ and $\chi_0$ the primitive character that induces it. A factor \m{1-\chi_0(p)p^{-s}} of
property~(d) vanishes only if \m{|\chi_0(p)|=p^{\re s}}; the left side is $0$ or~$1$, so this forces
\m{\re s=0}. Hence for \m{\re s>0}, and \m{s\ne1} when $\chi$ is principal, \m{L(s,\chi)=0} if and only if
\m{L(s,\chi_0)=0}. So Statement~3 is equivalent to Statement~4, and Statement~2 is equivalent to its restriction
to primitive characters. For a principal character, $\chi_0$ is the character modulo~$1$, whose $L$-function
is~$\zeta$: Statements 2 and~3 contain the corresponding statements for~$\zeta$. This remark is ours and is not a
declaration of the unit; the artifact proves only that Statement~3 implies Statement~4.
\end{remark}

\sect{Notes on the formal text}{s:notes}
\begin{enumerate}[label=\textup{\arabic*.},leftmargin=1.6em,itemsep=5pt]
\item \emph{The file.} The four definitions are the whole content of one file, whose only import is the library.
So the meaning of the four statements rests on the library alone, not on any definition of the artifact. This is
by design: it is the file to read in order to know what is claimed.
\item \emph{The modulus.} The formal text writes the modulus as a natural number~$N$ together with the
assumption \m{N\ne0}, which Lean passes as an instance argument; this is ``every integer \m{M\ge1}'' above. The
letter~$N$ is kept in this write-up for the level of the automorphic forms of the proof, and a modulus is
written~$M$ throughout, as in the weighted prime sums of the argument.
\item \emph{Statement 1 and the library's own formulation.} The library contains a formulation of the Riemann
hypothesis as a proposition. Statement~1 is that proposition, written out again in the artifact: the two texts
differ only in how the number \m{n+1} is read as a complex number. The first of the five final theorems has the
library's proposition as its statement and the second has Statement~1; the artifact proves the second by citing
the first, which Lean accepts only if the two propositions agree after unfolding.
\item \emph{The trivial zeros.} They appear as ``there is no natural number~$n$ with \m{s=-2(n+1)}'', that is,
\m{s\notin\{-2,-4,-6,\dots\}}.
\item \emph{The pole.} In Statement~1 the hypothesis \m{s\ne1} excludes the point where the value of~$\zeta$ is a
convention. In Statement~2 the hypothesis reads ``$\chi$ is not the unit element of the group of characters, or
\m{s\ne1}''; it excludes exactly the pole of a principal character. Statements 3 and~4 need no such hypothesis,
because \m{\re s<1} there.
\item \emph{The strip.} In Statements 3 and~4 the strip is open, and the two inequalities \m{0<\re s} and
\m{\re s<1} are two separate hypotheses, in this order, before the hypothesis \m{L(s,\chi)=0}.
\item \emph{Order of the quantifiers.} Statement~4 takes the modulus, the character and the hypothesis that the
character is primitive before the complex number~$s$; Statements 2 and~3 take the modulus, the character and
then~$s$. Nothing depends on this order here.
\item \emph{Real numbers.} The real part is a real number, and \m{\tfrac12} is the real number one half.
\end{enumerate}

\subhead{How this note was written.} From the dossier of the unit, which has no structured proof because the unit
has no proof; the program check of the dossier covered \UnitSteps{} leaf steps of \UnitProofs{} proofs and found
the four statements at their lines. I read the whole file of the unit in the artifact and compared each of the
four statements above with it, clause by clause. For the objects of the library I read, in the library's public
repository at the commit that the artifact pins, the declarations listed in the second table of Appendix~A: their
statements and the comments above them, not their proofs, and not the construction of the Hurwitz zeta functions
from which $\zeta$ and the $L$-functions are made. The dossier describes the same objects from memory and says so;
its descriptions agree with what I read. The two remarks of §\ref{s:proofs} are mine.

\head{Appendix A. Concordance with the Lean artifact}
{\fontsize{8.6}{11.6}\selectfont\renewcommand{\leanfont}{\ttfamily\fontsize{7.7}{10}\selectfont\color{leancol}}
\setlength{\LTleft}{0pt}\setlength{\LTright}{\fill}\setlength{\tabcolsep}{0pt}
\begin{longtable}{@{}p{30mm}p{56mm}p{43mm}r@{}}
\textbf{In this note} & \textbf{Lean declaration} & \textbf{File} & \textbf{Lines}\\\midrule\endhead
\ConcordanceRows
\end{longtable}}

\subhead{The four statements as they are written in the artifact.} Statement~1:
\quoted{rh}
Statement~2:
\quoted{halfPlane}
Statement~3:
\quoted{strip}
Statement~4:
\quoted{stripPrim}
In these texts the modulus is~$N$, the letter {\leanfont ℂ} stands for the complex numbers and {\leanfont ℕ} for
the natural numbers, {\leanfont s.re} is \m{\re s}, {\leanfont χ ≠ 1} says that $\chi$ is not the principal
character, and {\leanfont [NeZero N]} is the assumption \m{N\ne0}.

\subhead{The objects of the library.} The library is Mathlib at the commit that the artifact pins in its manifest.
The declarations below were read there for this note; the left column says what each one is in §\ref{s:setting}.

{\fontsize{8.6}{11.6}\selectfont\renewcommand{\leanfont}{\ttfamily\fontsize{7.7}{10}\selectfont\color{leancol}}
\setlength{\LTleft}{0pt}\setlength{\LTright}{\fill}\setlength{\tabcolsep}{0pt}
\begin{longtable}{@{}p{50mm}@{\hspace{4mm}}>{\raggedright\arraybackslash}p{86mm}@{}}
\textbf{In this note} & \textbf{Declaration of the library}\\\midrule\endhead
\mrow{the function $\zeta$}{riemannZeta}
\mrow{$\zeta$, property (a)}{zeta\_eq\_tsum\_one\_div\_nat\_cpow}
\mrow{$\zeta$, property (b)}{differentiableAt\_riemannZeta, riemannZeta\_residue\_one}
\mrow{$\zeta$, property (c)}{riemannZeta\_neg\_two\_mul\_nat\_add\_one, riemannZeta\_zero}
\mrow{the library's Riemann hypothesis}{RiemannHypothesis}
\mrow{Dirichlet character modulo $M$}{DirichletCharacter, MulChar}
\mrow{induced character}{DirichletCharacter.changeLevel}
\mrow{conductor, primitive character}{DirichletCharacter.conductor, DirichletCharacter.IsPrimitive}
\mrow{the function $L(s,\chi)$}{DirichletCharacter.LFunction}
\mrow{$L$, property (a)}{DirichletCharacter.LFunction\_eq\_LSeries}
\mrow{$L$, property (b)}{DirichletCharacter.differentiableAt\_LFunction}
\mrow{$L$, property (c)}{DirichletCharacter.LFunction\_modOne\_eq}
\mrow{$L$, property (d)}{DirichletCharacter.LFunction\_changeLevel}
\end{longtable}}

\head{Notation}
\begingroup\fontsize{8.6}{10.9}\selectfont\raggedright
\setlength{\columnsep}{9mm}\setlength{\parskip}{0pt}
\begin{multicols}{2}
\nt{$s$, $\re s$}{a complex number; its real part}
\nt{$n$}{a natural number, \m{n\ge0}}
\nt{$\zeta$}{the Riemann zeta function, as the library defines it}
\nt{$M$, $M_0$}{moduli of Dirichlet characters, integers \m{\ge1}; $M_0$ a divisor of~$M$}
\nt{$\chi$, $\chi_0$}{Dirichlet characters modulo $M$ and modulo~$M_0$}
\nt{$u$, $v$}{classes in \m{\ZZ/M}}
\nt{$L(s,\chi)$}{the $L$-function of the Dirichlet character~$\chi$, as the library defines it}
\nt{$p$}{a prime number}
\nt{$N$}{the modulus in the formal text; in this note it is~$M$}
\end{multicols}
\endgroup

\end{document}
